
\section{Experiments}

\subsection{Experimental Design}

The design is a full factorial of three primary demand levels (720, 1200 and
1680\,veh/h), four restriction configurations (none, upper, central or lower corridor)
and two signal plans (balanced and arterial priority), giving 24 contexts. Each is
simulated under six policy settings with three paired seeds, so the campaign consists of
$24\times6\times3=432$ runs. Demand levels are separated by 480\,veh/h, and that spacing
determines what the experiment can resolve about the location of a preference change
(Section~\ref{sec:boundary-exp}).

Four experiments are run on this campaign, ordered so that no later result can redefine an
earlier procedure. \textbf{E1} measures the switching margin between SP and DTT in every
context and determines, per restriction and signal-plan cell, the demand interval in which
it changes sign; \textbf{E2} compares the fixed policies, the ladder R0--R4 and the
hindsight bound using context-level regret; \textbf{E3} withholds entire demand levels to
separate interpolation inside the fitted support from extrapolation beyond it; and
\textbf{E4} runs the sensitivity checks and the two negative controls.

\subsection{Main Simulation Campaign}

All 432 runs completed and were retained; no run was discarded, re-run or replaced, and no
cell is missing. Accounting validity was checked before any decision analysis: the largest
discrepancy between edge-accumulated and trip-accumulated emission mass is
$1.69\times10^{-3}$ against a $10^{-2}$ tolerance, and the largest relative error of the
native emission algebra is $6.9\times10^{-7}$ against a $10^{-4}$ tolerance. A separate
route-probe campaign of 232 runs --- four contexts, two seeds and all 29 admissible paths
driven explicitly --- screened the emission forecast that the ECO and TECO policies
consume; it is the source of the decision to exclude ECO and is reported with the negative
controls in Section~\ref{sec:robust}.

\subsection{Switching-Boundary Resolution}
\label{sec:boundary-exp}

The experiment measures the switching margin at three demand levels per cell. When the
margin changes sign between adjacent levels it identifies the \emph{interval} containing
the change, not a threshold inside that interval, and no interpolated threshold is
reported anywhere in this paper. What the completed data can establish is whether each
bracket is itself resolved --- whether the observed sign of the margin at each endpoint is
large relative to dispersion across paired replications --- and that analysis, in
Section~\ref{sec:noise}, requires no additional simulation.

Localising the change would require simulating intermediate demand levels on this same
network: two additional levels across the eight cells and three portfolio policies, at
three seeds, is 144 runs. Those runs were not executed. The network and scenario
configuration files for this benchmark are not present in the analysis environment, and a
reconstructed network would not be comparable with the 432 completed runs, so the
additional runs could not be produced on the same footing as the existing evidence. No
claim here depends on them: every statement about the switching location is a statement
about an interval of the tested demand grid.

\subsection{Sensitivity and Robustness Experiments}

Fifteen combinations of tree depth $\in\{1,2,3,4,\text{unlimited}\}$ and minimum leaf size
$\in\{1,2,3\}$ are evaluated under the primary protocol; the time weight is swept over
$\{0,0.25,0.5,0.75,1\}$; the three seed blocks are each held out in turn; the scoring
scales are recomputed from training contexts only; and every decision is recomputed with
journey time replaced by in-network time, with the tail criterion in place of the mean,
and with route distance and then network stopped vehicle-time added as a third criterion.
The two negative controls are the ECO route-ordering gate and the weight sweep, both
designed to fail visibly if the corresponding component is doing no work.

\subsection{Statistical Treatment}

Each context is replicated three times with paired seeds, so a comparison between two
policies in one context is a paired difference $d_j$, $j=1,2,3$. We report the mean margin
$\bar d$ and the standard error $\mathrm{SE}(\bar d)=s_d/\sqrt{3}$, and call a margin
\emph{resolved} when $|\bar d|>2\,\mathrm{SE}(\bar d)$. This criterion is descriptive: it
marks which observed orderings are larger than the replication dispersion that produced
them. It is \textbf{not} a test of statistical significance --- with three replications the
standard error rests on two degrees of freedom, no distributional assumption is made, no
null hypothesis is specified and no multiplicity correction is applied. The paper contains
no hypothesis test and makes no claim of statistical significance anywhere. Aggregate
quantities such as mean regret are reported as benchmark statistics, and where their
sensitivity to replication matters the seed-block analysis of Section~\ref{sec:robust}
shows the spread directly rather than substituting an interval for it.
